\chapter{Tree-Based Models and Ensembles}
\label{chap:trees}
A decision tree models a prediction through successive conditional rules.
Tree ensembles then address different weaknesses of individual trees: bagging reduces instability by averaging diverse fits, whereas boosting adds functions that improve a loss.
This section develops their data flows, inductive biases, design choices and research trade-offs.
For the broader role of inductive bias, see \cref{sec:nfl-inductive-bias-the-assumptions-that-make-learning-possible}; for the compact model-family comparison, see \cref{sec:models-decision-trees-bagging-and-boosting}.

\section{A single decision tree: rules become regions}
\label{sec:trees-a-single-decision-tree-rules-become-regions}
Consider predicting machine-failure risk from temperature $T$ and vibration $V$.
A tree might apply ``$V>6$?'', then (if not) ``$T>70$?'', with a class or risk estimate at the end of each path.
Inference consists of traversing one path from the root to a \emph{leaf}; no optimisation happens at inference time.
The thresholds are learned during training, not specified by an expert in the standard algorithm.

\begin{figure}[htbp]\centering
\begin{tikzpicture}[node distance=9mm and 12mm,scale=.93,transform shape]
\node[flow] (root) {$V>6$?};
\node[yesbox,below left=of root] (high1) {High risk};
\node[flow,below right=of root] (temp) {$T>70$?};
\node[yesbox,below left=of temp] (high2) {High risk};
\node[yesbox,below right=of temp] (low) {Low risk};
\draw[arr] (root) -- node[above left,font=\scriptsize]{yes} (high1);
\draw[arr] (root) -- node[above right,font=\scriptsize]{no} (temp);
\draw[arr] (temp) -- node[above left,font=\scriptsize]{yes} (high2);
\draw[arr] (temp) -- node[above right,font=\scriptsize]{no} (low);
\end{tikzpicture}
\hfill
\begin{tikzpicture}[x=.075cm,y=.37cm,baseline=(current bounding box.center)]
\fill[palerteal] (0,0) rectangle (100,10);
\fill[red!17] (0,6) rectangle (100,10);
\fill[red!17] (70,0) rectangle (100,6);
\draw[->,thick] (0,0)--(107,0) node[right,font=\scriptsize] {$T$};
\draw[->,thick] (0,0)--(0,11.5) node[above,font=\scriptsize] {$V$};
\draw[dashed,amber,thick] (0,6)--(100,6);
\draw[dashed,amber,thick] (70,0)--(70,6);
\foreach \x in {0,70,100}{\draw (\x,0)--(\x,-.24) node[below,font=\tiny]{\x};}
\draw (0,6)--(-2,6) node[left,font=\tiny]{6};
\node[font=\scriptsize,align=center,text=navy] at (29,3) {Low\\risk};
\node[font=\scriptsize,align=center,text=navy] at (85,3) {High\\risk};
\node[font=\scriptsize,text=navy] at (45,8) {High risk};
\end{tikzpicture}
\caption{The same decision tree as a rule hierarchy (top) and a feature-space partition (bottom).
Every rectangular leaf region corresponds to one terminal prediction.}
\label{fig:tree_regions}
\end{figure}

A conventional \term{axis-aligned} tree splits on one feature at a time, $x_j\leq\theta_{\text{threshold}}$ versus $x_j>\theta_{\text{threshold}}$.
Its leaves divide feature space into $N_{\text{leaves}}$ disjoint regions.
A standard constant-leaf regression tree is
\begin{equation}
\boxed{
\hat f(x)
=
\sum_{j=1}^{N_{\text{leaves}}}
\underbrace{\theta_{\text{leaf},j}}_{\text{constant prediction in leaf }j}
\mathbb 1[x\in R_j]
}.
\end{equation}
Under squared-error loss, each $\theta_{\text{leaf},j}$ is usually the mean training target of the observations routed to that leaf.
A classification leaf can return the most frequent class or a vector of empirical class proportions.
Small-leaf proportions are not necessarily calibrated probabilities.
A tree therefore models a \emph{piecewise-constant function}, with conditional feature interactions created by its hierarchy: in \cref{fig:tree_regions}, temperature matters only for the branch with $V\leq6$.

\subsection{Learning a split}
\label{subsec:trees-learning-a-split}
At a node, a greedy learner generates allowable feature--threshold candidates, measures the reduction in a node criterion, selects the candidate with greatest gain, then recursively repeats on its children until a stopping rule applies.
The training process is a sequence of discrete structural decisions, rather than ordinary gradient descent through split thresholds.

\begin{figure}[htbp]\centering
\resizebox{\linewidth}{!}{%
\begin{tikzpicture}[node distance=6mm and 6mm]
\node[flow] (data) {Data reaching\\current node};
\node[flow,right=of data] (search) {Propose feature\\and threshold};
\node[flow,right=of search] (score) {Measure weighted\\impurity reduction};
\node[flow,right=of score] (pick) {Choose best\\allowed split};
\node[yesbox,right=of pick] (recur) {Route observations\\and repeat};
\draw[arr] (data)--(search);\draw[arr] (search)--(score);\draw[arr] (score)--(pick);\draw[arr] (pick)--(recur);
\end{tikzpicture}}
\caption{Recursive tree fitting: a local split is chosen before the resulting subproblems are solved.
A stopping rule instead produces a leaf prediction.}
\label{fig:tree_fit_flow}
\end{figure}

For $N_{\text{classes}}$ classes, a common impurity is \term{Gini impurity}, $G=1-\sum_{j=1}^{N_{\text{classes}}}p_{\text{class }j}^{2}$.
Gini is zero in a pure node; in balanced binary classification it is $0.5$.
Suppose eight examples have four positives and four negatives, so $G_{\rm parent}=0.5$.
A candidate split creates a five-sample child with class counts $(4,1)$ and a three-sample pure child.
The five-sample impurity is $1-(4/5)^2-(1/5)^2=0.32$, giving
\begin{equation}
\boxed{\Delta G=0.5-\left(\frac58\,0.32+\frac38\,0\right)=0.30.}
\end{equation}
A split is useful because it creates more concentrated class distributions.
\emph{Entropy} $-\sum_j p_{\text{class }j}\log p_{\text{class }j}$ is another classification criterion.
Regression trees usually select splits that reduce within-node squared error/variance, or another specified regression objective; MAE-based methods naturally use median leaf predictions.
These are alternative definitions of \emph{good partition}, not unrelated model families.
See the loss-function discussion in \cref{sec:objectives-supervised-losses-for-continuous-targets,sec:objectives-classification-losses-and-logits}.

\begin{calloutintuition}
A tree learns a conditional partition of feature space.
Standard leaves output constants; complex global predictions arise by routing inputs to different local regions.
\end{calloutintuition}

\section{Tree structure, greedy search and regularisation}
\label{sec:trees-tree-structure-greedy-search-and-regularisation}
\subsection{Greedy splits need not find the best tree}
\label{subsec:trees-greedy-splits-need-not-find-the-best-tree}
Local optimisation can miss feature interactions.
For balanced XOR data, $y=x_1\oplus x_2$, each binary input feature alone leaves a 50--50 class mixture in both children.
The first split achieves zero Gini gain, even though a depth-two tree can represent XOR exactly after splitting on both features.
A greedy implementation that insists on positive immediate gain may never find it.
This separates \emph{representational capacity} from \emph{optimisation ability}.

\begin{figure}[htbp]\centering
\begin{tikzpicture}[x=1.05cm,y=1.05cm]
\fill[red!15] (0,0) rectangle (1,1);\fill[palerteal] (1,0) rectangle (2,1);
\fill[palerteal] (0,1) rectangle (1,2);\fill[red!15] (1,1) rectangle (2,2);
\draw[thick] (0,0) rectangle (2,2);\draw[white,thick] (1,0)--(1,2) (0,1)--(2,1);
\node at (.5,.5){$0$};\node at (1.5,.5){$1$};\node at (.5,1.5){$1$};\node at (1.5,1.5){$0$};
\node[below] at (.5,0){$x_1=0$};\node[below] at (1.5,0){$x_1=1$};
\node[left] at (0,.5){$x_2=0$};\node[left] at (0,1.5){$x_2=1$};
\end{tikzpicture}
\caption{Balanced XOR: neither first-level split helps on its own, but two levels perfectly separate the classes.}
\label{fig:tree_xor}
\end{figure}

\subsection{Depth, leaves and pruning}
\label{subsec:trees-depth-leaves-and-pruning}
A binary tree with maximum depth $N_{\text{depth}}$ has at most $2^{N_{\text{depth}}}$ leaves; depth 1 is a \emph{decision stump}.
More leaves can approximate increasingly local structure, but individual leaf estimates may rely on few observations.
This is the bias--variance trade-off already introduced in \cref{sec:evaluation-overfitting-underfitting-and-the-bias-variance-tradeoff}.

\begin{center}\small
\begin{tabularx}{\linewidth}{@{}L{.24\linewidth}YY@{}}\toprule
Control & What it constrains & Trade-off\\\midrule
Maximum depth & Number of sequential decisions & Less depth limits interactions; greater depth can fit noise\\
Minimum samples per leaf & Evidence supporting each local estimate & Larger leaves stabilise predictions but smooth local structure\\
Minimum samples to split & Eligibility to expand a node & Earlier stopping versus finer partitions\\
Maximum leaf count & Number of distinct terminal regions & Direct budget for partition complexity\\
Minimum split improvement & Gain required to add a split & Reject weak gains, possibly including useful precursors\\
Pruning strength & Number of retained branches after growing & Simpler tree versus potentially discarded genuine structure\\\bottomrule
\end{tabularx}
\end{center}

\term{Pre-pruning} prevents a split when a limit is reached; \term{post-pruning} removes branches from a larger fitted tree.
Cost-complexity pruning makes the trade-off explicit:
\begin{equation}
\underbrace{J(T)}_{\text{selection objective}}
=
\underbrace{L_{\mathrm{train}}(T)}_{\text{fit}}
+
\underbrace{\alpha N_{\text{leaves}}(T)}_{\text{complexity cost}}.
\end{equation}
Each extra leaf must therefore buy enough training improvement to justify its penalty.
Pruning is not a complete solution to greedy-search errors: it simplifies already-discovered structures, rather than exploring every possible tree.

\subsection{Growth strategy and split geometry}
\label{subsec:trees-growth-strategy-and-split-geometry}
\term{Level-wise} growth expands eligible nodes layer by layer.
\term{Best-first/leaf-wise} growth expands the current leaf promising the greatest loss reduction.
Leaf-wise growth concentrates complexity in difficult regions but can create deep low-support paths; level-wise growth allocates depth more evenly.
Level-wise does \emph{not} imply symmetric splitting: nodes at the same depth can still use different tests.

An axis-aligned split tests one feature, such as $x_j>\theta_{\text{threshold}}$; an \term{oblique split} tests $\theta_{\text{split}}^\top x>\theta_{\text{threshold}}$.
A diagonal boundary, such as $T+2V>80$, may require many staircase-like axis-aligned regions but only one oblique split.
Oblique splits increase geometric expressivity and typically make split search harder.

\begin{calloutbias}
If the target is well described by hierarchical single-feature thresholds and locally constant regions, standard trees represent the structure efficiently.
If the boundary is smooth or strongly rotated, they may need many leaves; deeper trees add expressivity while increasing estimation variance.
\end{calloutbias}
Depth and leaf support are related but distinct controls.
\begin{calloutmemoryhook}
Depth limits interaction length; minimum leaf size limits evidence per local prediction.
\end{calloutmemoryhook}

\section{Bagging, Random Forests and ExtraTrees}
\label{sec:trees-bagging-random-forests-and-extratrees}
A single tree can be unstable: removing or changing a few observations may change a high-level split and many downstream predictions.
\term{Bagging} (bootstrap aggregating) reduces this variation by fitting multiple trees to bootstrap samples and aggregating their outputs.

\begin{figure}[htbp]\centering
\resizebox{.9\linewidth}{!}{%
\begin{tikzpicture}[node distance=7mm and 13mm]
\node[flow] (d) {Original data $D$};
\node[flow,below left=12mm and 35mm of d] (s1) {Bootstrap $D_1$};
\node[flow,below=12mm of d] (s2) {Bootstrap $D_2$};
\node[flow,below right=12mm and 35mm of d] (s3) {Bootstrap $D_3$};
\node[yesbox,below=of s1] (t1) {Tree 1};
\node[yesbox,below=of s2] (t2) {Tree 2};
\node[yesbox,below=of s3] (t3) {Tree 3};
\node[flow,below=21mm of t2] (agg) {Aggregate predictions};
\foreach \s in {s1,s2,s3}{\draw[arr] (d)--(\s);}
\foreach \i in {1,2,3}{\draw[arr] (s\i)--(t\i);\draw[arr] (t\i)--(agg);}
\end{tikzpicture}}
\caption{Bagging: independent bootstrap-fitted trees are combined at prediction time.
Individual fits can be trained in parallel.}
\label{fig:tree_bagging}
\end{figure}

Each bootstrap sample draws $N_{\text{examples}}$ observations \emph{with replacement} from $N_{\text{examples}}$ originals.
For large $N_{\text{examples}}$, about $63.2\%$ of originals appear at least once and about $36.8\%$ are out-of-bag for that tree.
Averaging regression predictions or class probabilities reduces variation when the trees' errors are not perfectly correlated.
Majority voting is another classification aggregation rule.

For $N_{\text{trees}}$ trees with common prediction variance $\sigma^2$ and pairwise correlation $\rho$, the equal-correlation model gives
\begin{equation}
\boxed{
\Var(\bar f)
=
\underbrace{
\underbrace{\sigma^2}_{\text{single-tree variance}}
\underbrace{\rho}_{\text{pairwise error correlation}}
}_{\text{shared variance: averaging cannot remove}}
+
\underbrace{
\sigma^2\frac{1-\rho}{N_{\text{trees}}}
}_{\text{unshared variance: shrinks with more trees}}
}.
\end{equation}
More trees only reduce the unshared component; decorrelation is therefore as important as tree count.
This is a simplified equal-variance/equal-correlation model, not a universal forest error law; connect it to the broader bias--variance picture in \cref{sec:evaluation-overfitting-underfitting-and-the-bias-variance-tradeoff}.

\term{Random Forests} add random selection of eligible features \emph{at each split}, commonly together with bootstrapping.
This decorrelates trees that might otherwise make the same dominant early choices.
Too much randomisation weakens individual trees, so diversity and individual predictive quality trade off.
\term{ExtraTrees} also randomise candidate thresholds and often use the original sample rather than bootstrapping by default.
The extra randomness can improve diversity and split-search efficiency while increasing individual-tree bias.

\begin{center}\small
\begin{tabularx}{\linewidth}{@{}L{.20\linewidth}YY@{}}\toprule
Method & Randomness & Primary idea\\\midrule
Bagged trees & Resample training rows & Average unstable learners\\
Random Forest & Resample rows; random features per node & Reduce correlation between learners\\
ExtraTrees & Random candidate thresholds; often random features & Additional split diversity, often less split-search work\\\bottomrule
\end{tabularx}
\end{center}

\term{Out-of-bag (OOB) evaluation} predicts each training observation using only trees whose bootstrap samples omitted it.
It offers an internal performance estimate under suitable sampling assumptions, but does not fix temporal or grouped leakage (\cref{sec:evaluation-cross-validation-and-fold-choice,sec:evaluation-model-selection-and-data-leakage}).

\begin{calloutmemoryhook}
Random Forests win by averaging diversity; more trees cannot remove shared error.
\end{calloutmemoryhook}

\section{Boosting: learning successive corrections}
\label{sec:trees-boosting-learning-successive-corrections}
Where bagging trains diverse trees independently, \term{boosting} trains a sequence of weak or limited-complexity learners that collectively improve a predictor.
\term{AdaBoost} reweights training examples according to previous classification errors; hard examples gain relative emphasis.
This can be useful, but mislabeled examples are also hard and may attract excessive attention.

\subsection{Gradient boosting as functional optimisation}
\label{subsec:trees-gradient-boosting-as-functional-optimisation}
\term{Gradient boosting} updates the current predictor by adding a new correction tree:
\begin{equation}
\boxed{
\underbrace{F_t(x)}_{\text{updated ensemble}}
=
\underbrace{F_{t-1}(x)}_{\text{current ensemble}}
+
\underbrace{\eta}_{\text{shrinkage}}
\underbrace{h_t(x)}_{\text{new correction tree}}
}.
\end{equation}
Each stage improves the current ensemble rather than refitting the original target from scratch.

For squared-error regression, suppose the current model predicts $5$ for targets $(3,5,7)$.
The residuals are $(-2,0,2)$.
If a correction tree fits these residuals exactly and $\eta=0.5$, the updated predictions become $(4,5,6)$.
The ensemble improves through a series of moderate adjustments.

More generally, each training example supplies the prediction-space direction
\begin{equation}
\boxed{
\underbrace{r_{i,t}}_{\text{pseudo-residual}}
=
-\left.\underbrace{\frac{\partial L(y_i,F)}{\partial F}}_{\text{loss slope in prediction space}}\right|_{F=F_{t-1}(x_i)}
}.
\end{equation}
A tree approximates this direction as a function of inputs; its terminal output values may subsequently be optimised for the specified loss.
Squared-error pseudo-residuals are ordinary residuals (up to loss convention).
For binary logistic loss differentiated with respect to logits, the negative gradient is label minus predicted probability.
This connects tree boosting to loss design in \cref{sec:objectives-supervised-losses-for-continuous-targets,sec:objectives-classification-losses-and-logits}.

\begin{figure}[htbp]\centering
\resizebox{\linewidth}{!}{%
\begin{tikzpicture}[node distance=6mm and 7mm]
\node[flow] (f) {Current model $F_{t-1}$};
\node[flow,right=of f] (p) {Predict training\\observations};
\node[flow,right=of p] (g) {Compute negative\\loss gradients};
\node[flow,right=of g] (t) {Fit correction\\tree $h_t$};
\node[yesbox,right=of t] (u) {Update\\$F_t=F_{t-1}+\eta h_t$};
\draw[arr] (f)--(p);\draw[arr] (p)--(g);\draw[arr] (g)--(t);\draw[arr] (t)--(u);
\draw[arr,rounded corners] (u.south)--++(0,-.7)--([yshift=-.7cm]f.south)--(f.south);
\end{tikzpicture}}
\caption{Gradient boosting is stage-wise optimisation in prediction-function space: fit a correction, add it, recompute the next direction.}
\label{fig:tree_boost_flow}
\end{figure}

This is related to gradient descent but not identical to backpropagating through a fixed neural network.
Neural-network training typically updates continuously parameterised weights within an existing architecture.
Boosting adds a \emph{new fitted function} at each iteration; hard tree split decisions are not usually differentiated end to end.
The SGD and adaptive-optimiser discussion in \cref{sec:optimisation-batch-gradient-descent-vs-sgd-vs-mini-batch-sgd,sec:optimisation-adam-momentum-plus-adaptive-step-sizes} concerns a different parameterisation of the optimisation problem.


\subsection{Ranking-aware boosting: LambdaMART}
\label{subsec:trees-lambdamart}
A ranking metric such as NDCG changes only when item order changes, so it is awkward as a direct smooth training loss (\cref{sec:objectives-ranking-and-preference-objectives,sec:evaluation-retrieval-and-ranking-metrics}).
LambdaRank-style training instead creates pairwise correction signals and weights them by the ranking consequence of swapping the pair.
If reversing two items would greatly change NDCG, that ordering error receives a stronger update than a swap deep in the list with almost no metric effect.

\begin{figure}[htbp]
\centering
\begin{tikzpicture}[node distance=8mm and 10mm]
\node[flow] (pair) {Misordered\\item pair};
\node[flow,right=of pair] (delta) {Estimate\\$|\Delta\mathrm{NDCG}|$};
\node[decisionnode,right=of delta] (big) {large\\effect?};
\node[yesbox,above right=4mm and 10mm of big] (strong) {Strong pairwise\\correction};
\node[yesbox,below right=4mm and 10mm of big] (weak) {Weak pairwise\\correction};
\node[flow,right=18mm of big] (tree) {Fit next\\boosting tree};
\draw[arr] (pair)--(delta);
\draw[arr] (delta)--(big);
\draw[arr] (big)--node[above,font=\scriptsize]{yes}(strong);
\draw[arr] (big)--node[below,font=\scriptsize]{no}(weak);
\draw[arr] (strong)--(tree);
\draw[arr] (weak)--(tree);
\end{tikzpicture}
\caption{LambdaMART keeps the stage-wise tree-boosting mechanism but makes the correction signal sensitive to where a ranking mistake occurs.}
\end{figure}

LambdaMART combines these ranking-aware ``lambda'' updates with MART-style gradient-boosted regression trees.
The important conceptual change is not a new tree primitive: it is a new objective signal layered onto the boosting machinery already developed above.

\begin{calloutmemoryhook}
LambdaMART is boosting where high-impact ranking mistakes get stronger correction signals.
\end{calloutmemoryhook}

\subsection{Curvature and second-order updates}
\label{subsec:trees-curvature-and-second-order-updates}
Second-order tree boosting also uses curvature of the loss with respect to the current prediction.
For one constant-output leaf, a common local quadratic update is
\begin{equation}
\boxed{
\underbrace{\theta_{\text{leaf update}}}_{\text{constant correction in this leaf}}
=
-\frac{
\underbrace{\sum_{i\in\mathrm{leaf}}\frac{\partial L_i}{\partial F_i}}_{\text{accumulated loss slope}}
}{
\underbrace{\sum_{i\in\mathrm{leaf}}\frac{\partial^2L_i}{\partial F_i^2}}_{\text{prediction-space curvature}}
+
\underbrace{\lambda_{\text{leaf}}}_{\text{shrink large leaf updates}}
}
}.
\end{equation}
This is scalar curvature in prediction space, not the full parameter Hessian; compare \cref{subsec:la-hessian,sec:optimisation-newton-s-method-and-other-second-order-methods}.

\subsection{Learning rate, depth and early stopping}
\label{subsec:trees-learning-rate-depth-and-early-stopping}
Smaller $\eta$ gives more conservative corrections, normally requiring more boosting stages.
Large $\eta$ can fit quickly but may overstep or fit noise.
The \emph{learning rate and number of trees must be tuned together}; doubling tree count does not exactly undo halving $\eta$, because the intermediate ensemble and all subsequent correction targets change.

A depth-one stump contributes a function of one feature; an additive ensemble of such stumps cannot represent genuine cross-feature interactions on its raw score scale without richer components.
Depth-two and deeper trees can encode conditional interactions within each stage.
Validation-based early stopping selects the iteration before continuing improvement on the training loss ceases to help generalisation (\cref{sec:stability-early-stopping-stop-before-the-model-overfits,sec:evaluation-model-selection-and-data-leakage}).

\begin{center}\small
\begin{tabularx}{\linewidth}{@{}L{.25\linewidth}YY@{}}\toprule
Design & Bagging / Random Forest & Gradient boosting\\\midrule
Tree dependence & Trees independently fitted & Each stage depends on preceding model\\
Combination & Average or vote & Add weighted corrections\\
Main benefit & Stabilise high-variance predictors & Build fit through stage-wise loss reduction\\
Typical trees & Often deep & Often shallow or leaf-limited\\
Parallelism & Trees can be trained concurrently & Boosting stages are sequential\\
Important controls & Tree count, randomness, leaf size & Learning rate, stages, leaf size, regularisation\\\bottomrule
\end{tabularx}
\end{center}

\begin{calloutpractical}
For gradient boosting, monitor validation loss across rounds and use early stopping.
A falling training loss does not prove the next tree improves deployment performance.
\end{calloutpractical}

\section{XGBoost, LightGBM and CatBoost: what changed?}
\label{sec:trees-xgboost-lightgbm-and-catboost-what-changed}
The historical progression is best seen as multiple related developments rather than one mandatory chain.
The useful story is the accumulation of problems and fixes, not a claim that each method replaced the previous one.

\begin{center}
\small
\begin{tabularx}{\linewidth}{@{}L{.16\linewidth}L{.22\linewidth}YY@{}}
\toprule
\textbf{Step} & \textbf{What changed} & \textbf{Why it was added} & \textbf{Consequence / new trade-off} \\
\midrule
CART & Recursive threshold partitioning & Represent nonlinear conditional rules & Flexible but unstable single trees \\
Bagging / Random Forest & Average diversified trees & Reduce variance from unstable fits & More compute; averaging cannot remove shared error \\
AdaBoost / Gradient Boosting & Add learners sequentially & Correct current mistakes or descend a loss & Strong staged fit; stages become dependent \\
XGBoost & Second-order updates, explicit penalties, scalable split search & Regularise and accelerate boosting & More design knobs and systems complexity \\
LightGBM / CatBoost & Selective computation, categorical handling, constrained growth choices & Improve efficiency or data handling & Extra approximation and inductive-bias choices \\
\bottomrule
\end{tabularx}
\end{center}

The modern libraries below share a core boosting framework but emphasise different implementation choices.

\subsection{XGBoost: explicit regularisation and efficient optimisation}
\label{subsec:trees-xgboost-explicit-regularisation-and-efficient-optimisation}
XGBoost popularised a scalable second-order tree-boosting system with explicit structural penalties.
A representative objective separates predictive fit from structural/leaf penalties:
\begin{equation}
\underbrace{\sum_i L(y_i,F(x_i))}_{\text{predictive loss}}
+
\sum_{t=1}^{N_{\text{trees}}}
\left(
\underbrace{\gamma N_{\text{leaves},t}}_{\text{cost of adding leaves}}
+
\underbrace{
\frac{\lambda}{2}
\sum_{j=1}^{N_{\text{leaves},t}}\theta_{\text{leaf},t,j}^{2}
}_{\text{shrink large leaf outputs}}
\right).
\end{equation}
The two penalties regulate structural complexity and correction magnitude separately.
This is regularisation of tree structure/leaf outputs rather than a generic neural-network weight penalty; compare \cref{sec:stability-regularisation-l1-and-l2}.
Other design choices include row/column subsampling, missing-value-aware routing, approximate or histogram split search, and efficient memory/parallel processing.
Leaf-count penalties and leaf-weight shrinkage regulate different forms of complexity.

\subsection{Histogram split search}
\label{subsec:trees-histogram-split-search}
Searching every distinct numerical threshold can be expensive.
A histogram learner maps numerical values to a finite set of bins, accumulates gradient/curvature statistics within each bin, and searches boundaries between bins.
For example, $(1.1,1.2,1.5,2.8,3.0,3.1,4.9,5.0)$ may be quantised into three bins containing $(1.1,1.2,1.5)$, $(2.8,3.0,3.1)$ and $(4.9,5.0)$.

This reduces split-search and memory cost at the expense of threshold resolution.
The approximation need not worsen validation performance, but its quality and runtime depend on bin count and the distribution of feature values.
Histogram methods now appear in several GBDT libraries; they are not unique to LightGBM.

\subsection{LightGBM: selective use of computation}
\label{subsec:trees-lightgbm-selective-use-of-computation}
LightGBM is associated with several efficiency choices:
\begin{itemize}
\item \term{Leaf-wise growth}: preferentially expand the leaf with greatest gain; efficient use of a leaf budget, but potentially deep low-support branches.
\item \term{GOSS (Gradient-based One-Side Sampling)}: keep high-gradient observations, sample some low-gradient observations, then reweight as appropriate for split-gain estimation.
This saves work while retaining influential gradient information; it is an available sampling strategy, not necessarily enabled in every configuration.
\item \term{EFB (Exclusive Feature Bundling)}: bundle sparse features that are mutually or nearly mutually exclusive, reducing effective dimensionality.
Exclusivity means seldom nonzero together, not merely low correlation.
\item \term{Histogram-based processing}: quantise features and aggregate split statistics, reducing work relative to exhaustive threshold scans.
\end{itemize}
The common principle is to allocate compute to observations, features and regions where it is likely to matter.
These approximations can alter accuracy and regularisation, and must be tested under equal resource budgets.

\subsection{CatBoost: ordered categories and symmetric trees}
\label{subsec:trees-catboost-ordered-categories-and-symmetric-trees}
Categorical features create a representation problem.
One-hot encoding can greatly expand high-cardinality features.
Global target encoding compresses them by replacing each category with a target statistic, but can leak an observation's own label into its input.
For a category represented by one example, naive full-training target encoding can directly expose the answer.

CatBoost uses \term{ordered categorical statistics}: conceptually, an example's category statistic depends on earlier observations in a chosen ordering, not on its own target, together with smoothing and other implementation details.
\term{Ordered boosting} is a distinct, related training procedure designed to reduce prediction shift arising from reused target information in gradient estimation.
Neither idea removes the need for appropriate temporal and group validation; a random permutation is not a replacement for chronology.

CatBoost also characteristically uses \term{symmetric (oblivious) trees}, where all nodes at the same depth use the same test.
This constrains the hypothesis class and supports efficient routing.
Unlike merely \emph{level-wise} trees, symmetric trees cannot choose separate splits at different nodes of one level.
CatBoost also offers other growth policies.

\begin{center}\footnotesize
\begin{tabularx}{\linewidth}{@{}L{.17\linewidth}YYY@{}}\toprule
 & XGBoost & LightGBM & CatBoost\\\midrule
Emphasis & Scalable regularised boosting & Efficiency and selective allocation & Categories, prediction shift, symmetric trees\\
Split search & Exact/approximate/histogram options & Histogram-based & Quantised feature processing\\
Growth & Depth-wise or loss-guided variants & Leaf-wise characteristically & Symmetric by default, alternatives\\
Categories & Native handling or external encoding & Native handling or encoding & Ordered statistics and other transforms\\
Characteristic ideas & Second-order updates, sparse-aware handling, explicit penalties & GOSS, EFB, histogram subtraction & Ordered boosting, ordered category statistics\\\bottomrule
\end{tabularx}
\end{center}

\begin{calloutpractical}
Library names are not independent hypotheses: compare specific growth, categorical-handling, sampling and regularisation mechanisms under a fair validation and compute budget.
\end{calloutpractical}

\section{Research optionality: what can be changed?}
\label{sec:trees-research-optionality-what-can-be-changed}
A tree system has several distinct design layers: input representation, base-tree structure, ensemble mechanism, optimisation objective, sampling procedure and computational implementation.
Different modifications target different limitations.

\small
\begin{longtable}{@{}L{.21\linewidth}L{.35\linewidth}L{.35\linewidth}@{}}
\caption{Degrees of experimentation for tree-based methods.}\label{tab:tree_research_axes}\\
\toprule Design axis & Representative alternatives & Mechanism to measure\\\midrule\endfirsthead
\toprule Design axis & Representative alternatives & Mechanism to measure\\\midrule\endhead
Split geometry & Axis-aligned versus oblique & Fit to rotated boundaries; split count; compute\\
Split search & Greedy versus lookahead & Discover interactions with initially weak splits\\
Leaf outputs & Constant versus linear local models & Approximation of smooth local behaviour\\
Growth policy & Level-wise versus best-first & Leaf depths, support and loss reduction\\
Ensemble & Single tree, bagging, boosting & Stability versus progressive loss reduction\\
Loss & Squared, robust, logistic, quantile & Alignment to task and error costs\\
Sampling & Bootstrap, feature subsets, gradient-aware & Prediction correlation, training cost and variance\\
Regularisation & Pruning, leaf penalties, shrinkage & Complexity, noisy-fit sensitivity\\
Categories & One-hot, ordered statistics, native splits & High-cardinality generalisation and leakage\\
Split resolution & Exact values versus histogram bins & Speed--accuracy--memory trade-off\\
Uncertainty & Point, quantile, distributional predictions & Calibration and coverage under assumptions\\
\bottomrule\end{longtable}
\normalsize

A research improvement can come from a richer \emph{hypothesis class} (oblique splits or linear leaves), better \emph{optimisation} (lookahead or second-order updates), better \emph{regularisation} (leaf penalties or shrinkage), better \emph{representation} (categorical statistics), or better \emph{systems engineering} (histograms and memory layout).
These are different causal explanations of performance gains.

Three useful controlled experiments are: compare bagging with Random Forests while logging tree correlation and ensemble error; compare level-wise with leaf-wise growth under a matched leaf budget while measuring leaf support and depth; and compare first- with second-order boosting using matched objectives while reporting wall-clock cost \emph{as well as} improvement per round.

\begin{calloutdegreeofexperiment}
Vary split geometry/search, growth policy, leaf model, ensemble type, objective, row/feature sampling, regularisation, category handling, histogram resolution, uncertainty output, and compute budget.
\end{calloutdegreeofexperiment}

\section{Inductive bias and model selection}
\label{sec:trees-inductive-bias-and-model-selection}
A tree's built-in preferences are: useful single-feature thresholds, hierarchical interactions, locally constant predictions, and adaptive refinement of regions.
These assumptions can help with heterogeneous tabular features and abrupt conditional effects.
They can be inefficient for oblique boundaries, smooth globally varying functions and behaviour far outside the training support (\cref{sec:nfl-inductive-bias-the-assumptions-that-make-learning-possible}).

\subsection{Extrapolation and feature transformations}
\label{subsec:trees-extrapolation-and-feature-transformations}
Suppose the physical target is $y=2x$ and training data only cover $0\leq x\leq10$.
A linear model may extrapolate the same slope to $x=15$.
A standard regression tree routes every sufficiently large $x$ to its terminal rightmost leaf and returns its constant estimate.
A finite sum or average of such trees likewise stays constant beyond all its final thresholds along that coordinate, holding other features fixed.
More trees improve within-support resolution but do not automatically recover the outside-support physical law.

\begin{figure}[htbp]\centering
\begin{tikzpicture}
\begin{axis}[width=.79\linewidth,height=5.0cm,xmin=0,xmax=16,ymin=0,ymax=34,axis lines=left,xlabel={$x$},ylabel={predicted $y$},xtick={0,5,10,15},ytick={0,10,20,30},grid=both,grid style={rulegray!50},legend style={at={(.03,.97)},anchor=north west,font=\small,draw=none,fill=white}]
\addplot[draw=none,fill=palerteal] coordinates {(0,0)(10,0)(10,34)(0,34)} \closedcycle;
\addplot[very thick,deepblue,domain=0:16] {2*x};\addlegendentry{linear relationship}
\addplot[very thick,amber,const plot] coordinates {(0,2)(2,6)(4,10)(6,14)(8,18)(10,19)(16,19)};\addlegendentry{tree-like staircase}
\addplot[dashed,slate,thick] coordinates {(10,0)(10,34)};
\end{axis}
\end{tikzpicture}
\caption{Schematic extrapolation: a constant-leaf tree saturates once an input lies beyond the last training-supported split; a linear function can continue its slope.
Curves are illustrative, not fitted experimental results.}
\label{fig:tree_extrapolation}
\end{figure}

Because exact axis-aligned splits primarily depend on feature ordering, strictly monotonic transforms of a numerical feature preserve the representable partitions (after transforming thresholds), unlike a general rotation that mixes features.
Practical fitted models may still change owing to histogram binning, ties and numerical detail.
Standard scaling is often unnecessary for exact tree splits; it may still matter for models or preprocessing stages around the tree (\cref{sec:representation-scaling-standardisation-and-leakage}).

\begin{calloutpractical}
Known physical law plus unexplained residuals: investigate $\hat y(x)=f_{\mathrm{physics}}(x)+f_{\mathrm{trees}}(x)$.
The physical term can encode extrapolative structure; the learned correction still needs out-of-domain validation.
\end{calloutpractical}

\begin{center}\small
\begin{tabularx}{\linewidth}{@{}L{.22\linewidth}YY@{}}\toprule
Model family & Preferred structure & Notable limitation to test\\\midrule
Linear / additive model & Global linear or additive effects & Missing interactions and discontinuities\\
Single tree & Hierarchical threshold rules & Instability; staircase approximation\\
Random Forest & Average diverse local partitions & Correlated errors; extrapolation\\
Boosted trees & Sequential additive partition corrections & Overfit to residual noise; extrapolation\\
Neural network & Learned continuous representations & Data/optimisation and architectural bias\\
Physics-informed/hybrid & Governing laws plus fitted deviations & Error from incorrect scientific assumptions\\\bottomrule
\end{tabularx}
\end{center}

\section{Evaluation, explanation and failure modes}
\label{sec:trees-evaluation-explanation-and-failure-modes}
\subsection{Evaluate the deployment task, not a random row split}
\label{subsec:trees-evaluate-the-deployment-task-not-a-random-row-split}
In longitudinal data, measurements from one machine or person can be strongly dependent.
Randomly dividing rows may put the same entity in training and test sets.
For forecasting, training on later observations than those being tested can introduce future-information leakage.
Choose chronological and/or grouped validation according to the actual deployment question (\cref{sec:evaluation-cross-validation-and-fold-choice,sec:adaptation-domain-adaptation-the-data-distribution-changes}).

Supervised feature engineering is another leakage route.
Computing global category target means, performing label-dependent feature selection before splitting, or using a row's own label during its target encoding can all inflate apparent validation results.
Fit transformations within the training fold, with cross-fitting or ordered encodings when appropriate.
Missingness may itself be predictive, but changing collection practices can create shift; an integer category code should not silently become an ordinal physical quantity.

\subsection{Metrics, importance and causal interpretation}
\label{subsec:trees-metrics-importance-and-causal-interpretation}
Rare-event prediction should consider precision, recall and PR curves, not accuracy alone; log loss and calibration measure different properties from ranking AUC.
Operating thresholds should reflect false-alarm and missed-event costs.
See the worked definitions in \cref{sec:evaluation-precision-recall-and-f1,sec:evaluation-roc-curves-auc-roc-and-precision-recall,sec:evaluation-decision-thresholds}.

Three common explanations answer different questions:
\begin{itemize}
\item \term{Impurity importance} accumulates training-criterion reductions attributed to features; continuous/high-cardinality features can receive a structural advantage.
\item \term{Permutation importance} measures held-out performance deterioration when one feature is shuffled; correlated predictors can substitute for one another, masking individual importance.
\item \term{SHAP} assigns additive contributions to a prediction relative to a chosen reference/distributional convention; an attribution is not an intervention or an automatic causal effect.
\end{itemize}
Partial dependence averages predictions while changing a chosen feature; individual conditional expectation traces this for one observation.
With strongly dependent features, intervening on one coordinate while holding others unchanged can generate unrealistic feature combinations.
Such plots describe a fitted model, not necessarily the true data-generating mechanism.

\begin{center}\small
\begin{tabularx}{\linewidth}{@{}L{.27\linewidth}YY@{}}\toprule
Failure mode & Mechanism & Useful check\\\midrule
Overfitting & Tiny leaves; late boosting stages fit noise & Leaf limits, early stopping, held-out trends\\
Leakage & Targets or future/group information in features & Pipeline audit, grouped/time splits\\
Poor calibration & Predicted scores do not match event frequencies & Reliability curves; probability scores such as log loss or Brier score\\
Shift & Deployment population differs from training & Time/site holdouts; drift monitoring\\
Misleading importance & Correlated predictors or split-cardinality effects & Refit stability; grouped perturbations\\
Slow inference & Many trees, deep paths, expensive preprocessing & Profile end-to-end latency and memory\\
Extrapolation error & Terminal leaves return learned local values & Validate beyond training support; hybrid model\\\bottomrule
\end{tabularx}
\end{center}

\begin{calloutmemoryhook}
Feature attribution describes model dependence, not causal mechanism.
\end{calloutmemoryhook}

\section{Extensions and a complete research example}
\label{sec:trees-extensions-and-a-complete-research-example}
\subsection{Beyond standard trees}
\label{subsec:trees-beyond-standard-trees}
Conventional trees are not the end of the design space:
\begin{itemize}
\item \term{Differentiable trees} soften routing, for example $p_{\mathrm{right}}(x)=\sigma((\theta_{\text{split}}^\top x-\theta_{\text{threshold}})/\tau_{\text{routing softness}})$, enabling gradients through route probabilities.
Small $\tau_{\text{routing softness}}$ approaches hard routing but may cause saturated gradients; soft routing can also make inference costlier.
\item \term{Model trees} predict $\theta_{\text{leaf},j}^\top x+\theta_{\text{bias},j}$ inside region $R_j$ rather than a constant.
This fits smooth local effects more directly, but local parameter fitting and continuity across boundaries become design issues.
\item \term{Quantile or distributional boosting} targets uncertainty rather than only a conditional mean.
Conformal calibration can provide prediction sets when future cases behave like the calibration cases under the method's stated assumptions; coverage under shift still requires examination.
\item \term{Constrained trees} may enforce monotonic effects or limit interactions, adding domain-specific bias.
An incorrect constraint can create systematic error.
\item \term{Neural--tree hybrids and tabular foundation models} change representation learning or transfer information across tasks rather than only improving the split search of a newly fitted tree.
\end{itemize}

\subsection{Worked design: failure forecasting on repeated measurements}
\label{subsec:trees-worked-design-failure-forecasting-on-repeated-measurements}
Suppose 500 machines produce repeated sensor measurements over several years.
The inputs include numerical readings and categorical manufacturers; failures are rare; labels are noisy and inference must be fast.
A defensible workflow is:
\begin{enumerate}
\item Define the unit of prediction (next-month failure per machine) and the information actually available at decision time.
\item Use chronological validation, and hold out complete machines as well if the deployment goal includes unseen machines.
Repeated rows do not equal independent machines.
\item Fit a simple physical or parametric baseline, a regularised single tree, a Random Forest and a regularised GBDT; evaluate suitable categorical treatment without target leakage.
\item Compare minimum leaf size, tree complexity, boosting shrinkage and early stopping; report stability across machines, time windows and random seeds.
\item Evaluate PR behaviour, calibration and application-specific false-alarm costs, not accuracy in isolation.
\item Profile end-to-end inference latency and memory.
Investigate whether manufacturer effects and feature importances remain stable across valid holdouts.
\end{enumerate}
This experiment tests specific modelling assumptions as well as parameter values.
The appropriate ensemble cannot be inferred from its library name alone.

\begin{calloutmemoryhook}
Trees partition; bagging averages diverse partitions; boosting adds loss-driven corrections.
\end{calloutmemoryhook}

\section{Further reading}
\label{sec:trees-further-reading}
\begin{itemize}
\item Chen and Guestrin (2016), \href{https://arxiv.org/abs/1603.02754}{\emph{XGBoost: A Scalable Tree Boosting System}}: regularised boosting and scalable split search.
\item Ke et al. (2017), \href{https://papers.nips.cc/paper_files/paper/2017/hash/6449f44a102fde848669bdd9eb6b76fa-Abstract.html}{\emph{LightGBM: A Highly Efficient Gradient Boosting Decision Tree}}: GOSS and EFB.
\item Prokhorenkova et al. (2018), \href{https://arxiv.org/abs/1706.09516}{\emph{CatBoost: Unbiased Boosting with Categorical Features}}: ordered statistics and boosting.
\item Grinsztajn et al. (2022), \href{https://arxiv.org/abs/2207.08815}{\emph{Why Do Tree-Based Models Still Outperform Deep Learning on Typical Tabular Data?}}: comparative inductive-bias study with a limited benchmark scope.
\end{itemize}
